% Chapter 3, Figure 17: control surface for calculating the friction drag due to a turbulent boundary layer
% over a flat plate at zero angle of attack (scan: figures/ch3/fig17.png).
% fig17.py: boundary-layer edges delta ~ x^(4/5), eq. (80); the profile at the trailing edge O the 1/7th-power
% law, eq. (74); the wake profile u(x,y) at B B_1 a Gaussian defect with the same momentum thickness,
% eqs. (73), (76), reaching U_inf at the wake lines (+-delta); uniform U_inf at A A_1. Data: fig17-edge.csv,
% fig17-O.csv, fig17-B.csv, fig17-arrows.csv. The control surface's sides are hatched inside, as printed.
% Overlay on the scan (fig17.calib.json), 95% within: edges 12 and 16 px, profile at O 6 px, wake profile
% 11.5 px: the 1973 art draws laminar-looking edges (about x^0.4) and a deeper wake (corrections/v2-figures.md).
% Drawn in the scan's pixels, y up from the plate, 1.2 times the printed size.
\documentclass[11pt]{standalone}
\usepackage{tamrfig}
% the profile family's styles (Ch3 Figs 16, 17, 18, 25)
\tikzset{
  profile/.style={draw=s1, line width=1pt},
  profile arrow/.style={draw=s1, line width=0.5pt, -{Stealth[length=3.4pt, width=2.4pt]}},
  profile stroke/.style={draw=s1, line width=0.5pt},
  plate/.style={draw=ink, line width=1.6pt, line cap=butt},
}
\pgfplotstableread[col sep=comma]{figures/v2/ch3/fig17-arrows.csv}\arrows
\pgfplotsset{canvas/.style={hide axis, x=0.02032cm, y=0.02032cm, disabledatascaling, anchor=origin,
  at={(0,0)}, xmin=0, xmax=1, ymin=0, ymax=1, enlargelimits=false, clip=false}}
\begin{document}
\begin{tikzpicture}[x=0.02032cm, y=0.02032cm]
  \def\xA{32.5}\def\xLE{108}\def\xO{340}\def\xB{449}\def\top{127}\def\bot{-127}\def\dl{58}\def\U{56}
  % the control surface A A_1 B_1 B: hatched inside
  \fill[hatch] (\xA,\top-8) rectangle (\xB,\top);
  \fill[hatch] (\xB-8,0) rectangle (\xB,\top);
  \fill[hatch] (\xA,0) rectangle (\xB,8);
  \fill[hatch] (\xA,0) rectangle (\xA+8,\top);
  \draw[outline] (\xA,0) rectangle (\xB,\top);
  % boundary-layer edges to O, the wake lines beyond (interrupted by the profiles)
  \begin{axis}[canvas]
    \addplot[guide, mark=none] table[x=x, y=y, col sep=comma] {figures/v2/ch3/fig17-edge.csv};
    \addplot[guide, mark=none] table[x=x, y expr=-\thisrow{y}, col sep=comma] {figures/v2/ch3/fig17-edge.csv};
  \end{axis}
  \foreach \s in {1,-1} {
    \draw[guide] (\xO+\U,\s*\dl) -- (\xB,\s*\dl);
    \draw[draw=ink2, line width=0.45pt] (\xB+\U,\s*\dl) -- (560,\s*\dl);
  }
  % the plate, its trailing edge O
  \draw[plate] (\xLE,0) -- (\xO,0);
  % profiles: arrows, curves, base lines, the U_inf references inside the layer
  \pgfplotstablegetrowsof{\arrows}\pgfmathtruncatemacro\nrows{\pgfplotsretval-1}
  \foreach \r in {0,...,\nrows} {
    \pgfplotstablegetelem{\r}{x0}\of\arrows\let\ax\pgfplotsretval
    \pgfplotstablegetelem{\r}{y}\of\arrows\let\ay\pgfplotsretval
    \pgfplotstablegetelem{\r}{x1}\of\arrows\let\bx\pgfplotsretval
    \pgfplotstablegetelem{\r}{head}\of\arrows\let\ah\pgfplotsretval
    \ifnum\ah=1 \draw[profile arrow] (\ax,\ay) -- (\bx,\ay); \else \draw[profile stroke] (\ax,\ay) -- (\bx,\ay); \fi
  }
  \draw[profile] (\xA+\U,\bot) -- (\xA+\U,\top);
  \begin{axis}[canvas]
    \addplot[profile, mark=none] table[x=x, y=y, col sep=comma] {figures/v2/ch3/fig17-O.csv};
    \addplot[profile, mark=none] table[x=x, y=y, col sep=comma] {figures/v2/ch3/fig17-B.csv};
  \end{axis}
  \draw[outline] (\xA,\bot) -- (\xA,0) (\xO,\bot) -- (\xO,\top) (\xB,\bot) -- (\xB,0);
  \draw[guide] (\xO+\U,-\dl) -- (\xO+\U,\dl) (\xB+\U,-\dl) -- (\xB+\U,\dl);
  \node[open point] at (\xO,0) {};
  \foreach \x in {\xA, \xO, \xB} \node[anchor=north, inner sep=3pt] at (\x+\U/2,\bot) {$U_\infty$};
  % flow leaving through A_1 B_1
  \draw[thin vec] (117,113) .. controls (150,121) and (180,137) .. (213,141);
  \draw[thin vec] (230,113) .. controls (263,121) and (293,137) .. (326,141);
  % corner letters, O, the wall shear, the wake profile's name
  \node[anchor=east, inner sep=2pt] at (\xA,\top) {$A_1$};
  \node[anchor=east, inner sep=2pt] at (\xA,0) {$A$};
  \node[anchor=south, inner sep=2pt] at (\xB,\top) {$B_1$};
  \node[anchor=north east, inner sep=1.5pt] at (\xB,0) {$B$};
  \node[anchor=north east, inner sep=1.5pt] at (\xO-1,-1) {$O$};
  \draw[thin vec] (252,15) -- (306,15) node[midway, above=1pt] {$\tau_o$};
  \node[anchor=west, inner sep=2pt] at (\xB+\U+4,22) {$u(x,y)$};
  % the control surface's name
  \node[align=left, anchor=south west, inner sep=1.5pt] (cs) at (62,136) {control\\surface};
  \draw[leader] (cs.east) -- (150,124);
  % coordinates and dimensions
  \draw[thin vec] (\xLE,0) -- (\xLE,80) node[above=1pt] {$y$};
  \draw[extension] (\xLE,-3) -- (\xLE,-95);
  \draw[thin vec] (\xLE,-62) -- (168,-62) node[right=1pt] {$x$};
  \dimline{(\xLE,-87.5)}{(\xO,-87.5)}{$\ell$}
  \dimline[pos=0.68]{(423,0)}{(423,\top)}{$h$}
  \dimline{(423,0)}{(423,-\dl)}{$\delta$}
\end{tikzpicture}
\end{document}
